\section*{Bab \ref{ch:b2:isotope-geochemistry} --- Geokimia Isotop}

\begin{solution}{exo:b2:isotope-geochemistry:1}
Tristan: $(0.51249/0.512638 - 1)\times10^4 = -2.9$. Mangaia: $(0.51285/0.512638 -
1)\times10^4 = +4.1$. MORB: $+8.5$. Sumber Mangaia dan sumber MORB, dengan
$\varepsilon_{\mathrm{Nd}}$ positif, rata-rata pernah memiliki Sm/Nd lebih tinggi daripada \omterm{def:b2:planet-formation-meteorites:chondrule}{kondrit};
sumber Tristan lebih rendah.
\end{solution}

\begin{solution}{exo:b2:isotope-geochemistry:2}
$\delta^{18}\mathrm{O}_{\mathrm{VPDB}} = (\num{2.0672e-3}/\num{2.0052e-3} - 1)\times1000 =
+30.9$ permil pada skala \omterm{def:b2:isotope-geochemistry:standards}{VSMOW}. Kalsit itu memiliki $^{18}\mathrm{O}/^{16}\mathrm{O} =
\num{2.0672e-3}\times(1 - 0.0020) = \num{2.0631e-3}$, sehingga $\delta^{18}\mathrm{O} =
(\num{2.0631e-3}/\num{2.0052e-3} - 1)\times1000 = +28.9$ permil (\omterm{def:b2:isotope-geochemistry:standards}{VSMOW}).
\end{solution}

\begin{solution}{exo:b2:isotope-geochemistry:3}
Pada \qty{1.0}{Ga}: $e^{\lambda t} - 1 = 0.006545$, dan $0.512638 - 0.1966\times0.006545 =
0.511351$. Pada \qty{4.567}{Ga}: $e^{\lambda t}-1 = 0.030243$, dan $0.512638 -
0.1966\times0.030243 = 0.506692$.
\end{solution}

\begin{solution}{exo:b2:isotope-geochemistry:4}
$\delta\mathrm{D} = 8\times(-18.0) + 10 = -134$ permil. Untuk danau, $d = -42 -
8\times(-4.0) = -10$ permil, di bawah nilai 10 milik garis itu: air danau terletak di bawah \omterm{prop:b2:isotope-geochemistry:mwl}{garis air
meteorik}, seperti air permukaan yang menguap (isotop ringan pergi lebih dulu).
\end{solution}

\begin{solution}{exo:b2:isotope-geochemistry:5}
$Q = 10^4\times\num{6.524e-12}\times0.1966/0.512638 = \num{2.502e-8}$ per tahun, atau
\qty{25.0}{\per\Ga}. Granit itu memiliki $\mathrm{Sm}/\mathrm{Nd} = 6.0/30 = 0.200$,
dibandingkan $0.145/0.457 = 0.317$ untuk CI: $f = 0.200/0.317 - 1 = -0.37$. Selama
\qty{0.5}{Ga}: $\Delta\varepsilon = 25.0\times(-0.37)\times0.5 = -4.6$ satuan.
\end{solution}

\begin{solution}{exo:b2:isotope-geochemistry:6}
$\varepsilon \approx 3.5 + 25.0\times(-0.45)\times2.7 = 3.5 - 30.4 = -26.9$, dibandingkan
$-27.3$: pendekatan itu 0.4 satuan terlalu tinggi (kurang negatif). Sebagian besar
selisihnya, sekitar 0.3 satuan, berasal dari $e^{\lambda t} > 1$: induknya, dan karena itu
kekurangan pertumbuhan relatif terhadap CHUR, dahulu lebih melimpah daripada sekarang
(\cref{exo:b2:isotope-geochemistry:10}); sisanya berasal dari pembulatan $f$ dan nilai awal.
\end{solution}

\begin{solution}{exo:b2:isotope-geochemistry:7}
Dengan fraksi massa granit $x$, $\mathrm{Sr} = 320x + 130(1-x)$ dan
$R = (320x\times0.7200 + 130(1-x)\times0.7030)/\mathrm{Sr}$. Untuk $x = 0.05$: Sr
\qty{139.5}{ppm}, $R = 0.70495$; $x = 0.10$: 149.0, 0.70665; $x = 0.20$: 168.0,
0.70948. Dari bukti proposisi pencampuran, $R\,\mathrm{Sr}$ linear terhadap Sr:
$R = 0.731632 - 3.7221/\mathrm{Sr}$, garis lurus dengan kemiringan $-3.722$ (ppm) dan
titik potong 0.7316 terhadap $1/\mathrm{Sr}$. Pemeriksaan pada $x = 0.05$: $0.731632 - 3.7221/139.5
= 0.70495$.
\end{solution}

\begin{solution}{exo:b2:isotope-geochemistry:8}
Pada \qty{20}{\degreeCelsius}: $\delta^{18}\mathrm{O} = 1000/1.0098 - 1000 = -9.70$
permil, $\delta\mathrm{D} = 1000/1.084 - 1000 = -77.5$ permil, kemiringan $77.5/9.70 = 7.99$.
Pada \qty{0}{\degreeCelsius}: $-11.56$ dan $-99.9$ permil, kemiringan 8.64. Keduanya dekat dengan
kemiringan 8 \omterm{prop:b2:isotope-geochemistry:mwl}{garis air meteorik}: fraksinasi kesetimbangan air
antara cair dan uap saja sudah menetapkan perbandingan pergeseran deuterium dan oksigen-18
di sekitar 8.
\end{solution}

\begin{solution}{exo:b2:isotope-geochemistry:9}
Kerak memuat $0.0051\times27/1.25 = 0.110$, yaitu 11 persen neodimium. Neraca
massa: $0.110\times(-15) + 0.890\,\varepsilon_m = 0$, sehingga $\varepsilon_m = 1.65/0.890 =
+1.9$. Dengan reservoir terkuras yang memuat bagian $x$ pada $+8.5$ dan sisanya pada 0:
$0.110\times(-15) + 8.5x = 0$, $x = 1.65/8.5 = 0.19$: sekitar seperlima
neodimium.
\end{solution}

\begin{solution}{exo:b2:isotope-geochemistry:10}
Laju pertumbuhannya $\lambda p\,e^{\lambda t}$, dengan $e^{\lambda t} \approx 1 + \lambda t$;
dirata-ratakan atas selang dari pemisahan $T$ hingga sekarang, faktornya
$1 + \lambda T/2$, sehingga suku-suku yang diabaikan berukuran relatif $\lambda T/2$. Sepanjang
umur Bumi, $\num{6.524e-12}\times\num{4.567e9}/2 = 0.015$, yaitu 1.5 persen. Untuk
kerak pada \cref{exo:b2:isotope-geochemistry:6}, $\lambda T/2 = 0.0088$, dan galatnya $0.0088\times30.4 = 0.27$
satuan. Perubahan lambat perbandingan CHUR yang menjadi pembagi $\varepsilon$ menambahkan
suku lain yang lebih kecil, berorde $(p/R)\lambda T \times \varepsilon_0$, di sini 0.02 satuan.
\end{solution}

\begin{solution}{exo:b2:isotope-geochemistry:11}
$K = 13.3/5.0 = 2.66$. Dengan $s = 0.5$: $s' = 2.66\times0.5/(1 + 1.66\times0.5) =
1.33/1.83 = 0.73$. Selisih $s' - s$ terbesar ketika
$\mathrm{d}s'/\mathrm{d}s = K/(1+(K-1)s)^2 = 1$, yaitu pada $s = (\sqrt K - 1)/(K-1) =
0.631/1.66 = 0.38$, dengan $s' = 0.62$: selisih 0.24.
\end{solution}

\begin{solution}{exo:b2:isotope-geochemistry:12}
Pembentukan: $2.3$ hingga $2.6\times\qty{7.2e9}{km^3}/\qty{4.5e9}{yr} = 3.7$ hingga
\qty{4.2}{km^3} per tahun. Penghancuran: $1.3$ hingga $1.6\times7.2/4.5 = 2.1$ hingga
\qty{2.6}{km^3} per tahun. Laju saat ini berbeda: penambahan (2.7) di bawah pembentukan
rata-rata model, penyingkiran (3.2) di atas penghancuran rata-ratanya. Dalam model itu kerak tumbuh
pada masa awal, ketika pembentukan paling cepat, dan masa kini adalah masa seimbang atau sedikit
berkurang, seperti dikatakan neraca zona subduksi.
\end{solution}

\begin{solution}{pb:b2:isotope-geochemistry:1}
\textbf{1.} Mantel: $(0.51307/0.512638 - 1)\times10^4 = +8.4$. Sedimen:
$(0.51141/0.512638 - 1)\times10^4 = -24.0$.

\textbf{2.} Kora $+2.8$, Vel $-4.3$, Asha $+4.7$.

\textbf{3.} Dengan $^{87}\mathrm{Sr}/^{86}\mathrm{Sr}$ mendatar dan $\varepsilon_{\mathrm{Nd}}$ tegak:
mantel (0.70282, $+8.4$), Kora (0.70410, $+2.8$), Vel (0.70618, $-4.3$), sedimen
(0.72000, $-24.0$), Asha (0.70290, $+4.7$). Hanya Vel, dengan
$\varepsilon_{\mathrm{Nd}}$ negatif, yang terletak di sisi CHUR yang diperkaya.

\textbf{4.} Di bawah larik: pada perbandingan stronsium \omterm{def:b2:isotope-geochemistry:reservoirs}{mantel terkuras}, larik itu
akan memiliki $\varepsilon_{\mathrm{Nd}}$ di sekitar $+8$; Asha memiliki $+4.7$. Asha mirip dengan
kelompok St.~Helena tipe HIMU (Mangaia, Tubuai, St.~Helena).

\textbf{5.} Peleburan mengubah perbandingan induk--anak tetapi tidak, secara terukur,
perbandingan isotop stronsium dan neodimium (proposisi tentang peleburan); lava-lava itu
muda, sehingga perbandingannya belum bertambah sejak erupsi (pertanyaan 6).

\textbf{6.} $\Delta R = 0.12\,(e^{\lambda\cdot10^6} - 1) = 0.12\times\num{6.52e-6} =
\num{7.8e-7}$, atau $\num{7.8e-7}/0.512638\times10^4 = 0.015$ satuan epsilon: jauh di bawah
perbedaan antarpulau, dan di bawah ketelitian pengukuran.

\textbf{7.} $\mathrm{Nd} = 40f + 0.75(1-f)$, dan
$R = \bigl(40f\times0.51141 + 0.75(1-f)\times0.51307\bigr)/\mathrm{Nd}$.

\textbf{8.} $f = 0.5$ persen: Nd \qty{0.9463}{ppm}, $\varepsilon = +1.6$; 1 persen:
1.1425, $-2.9$; 2 persen: 1.535, $-8.4$; 5 persen: 2.7125, $-15.4$.

\textbf{9.} Sr 10.95, 11.90, 13.80, dan \qty{19.50}{ppm}; $^{87}\mathrm{Sr}/^{86}\mathrm{Sr} =
0.70439$, 0.70571, 0.70780, dan 0.71163.

\textbf{10.} $R\,\mathrm{Sr} = 10\times0.70282 + f(200\times0.72000 - 10\times0.70282)$ dan
$\mathrm{Sr} = 10 + 190f$; dengan mengeliminasi $f$, $R\,\mathrm{Sr} = 7.0282 +
(\mathrm{Sr}-10)\times136.9718/190$, sehingga $R = 0.720904 - 0.18084/\mathrm{Sr}$: garis
lurus dengan kemiringan $-0.1808$ (ppm) terhadap $1/\mathrm{Sr}$. Pemeriksaan pada $f = 1$ persen:
$0.720904 - 0.18084/11.90 = 0.70571$.

\textbf{11.} $K = (10/0.75)/(200/40) = 13.33/5.0 = 2.67$. Mantel relatif
lebih kaya stronsium, sedimen lebih kaya neodimium, sehingga perbandingan Nd bergerak menuju
sedimen lebih cepat: sedikit sedimen menggeser $\varepsilon_{\mathrm{Nd}}$ jauh sebelum
$^{87}\mathrm{Sr}/^{86}\mathrm{Sr}$ bergerak banyak.

\textbf{12.} 53 kali lebih kaya Nd ($40/0.75$), 20 kali lebih kaya Sr. Pada $f = 1$ persen
sedimen menyumbang $0.40/1.1425 = 35$ persen neodimium, yang menggeser
$\varepsilon_{\mathrm{Nd}}$ sebesar $0.35\times(8.4 + 24.0) = 11.3$ satuan, tetapi hanya $2.0/11.9 =
17$ persen stronsium.

\textbf{13.} $0.01\times40/1.1425 = 0.350$: 35 persen.

\textbf{14.} Dari $R\,(40f + 0.75(1-f)) = 40fR_S + 0.75(1-f)R_D$:
$f = 0.75(R - R_D)/\bigl(40(R_S - R) + 0.75(R - R_D)\bigr)$, dan hal yang sama dengan 200
dan 10 untuk stronsium.

\textbf{15.} $f = 0.75\times(-0.00029)/\bigl(40\times(-0.00137) + 0.75\times(-0.00029)\bigr)
= \num{-2.175e-4}/\num{-5.502e-2} = 0.0040$: 0.40 persen.

\textbf{16.} $f = 10\times0.00128/\bigl(200\times0.01590 + 10\times0.00128\bigr) =
0.0128/3.1928 = 0.0040$: 0.40 persen; keduanya sesuai.

\textbf{17.} Neodimium: $0.75\times(-0.00065)/\bigl(40\times(-0.00101) +
0.75\times(-0.00065)\bigr) = \num{-4.875e-4}/\num{-4.089e-2} = 0.0119$; stronsium:
$10\times0.00336/\bigl(200\times0.01382 + 10\times0.00336\bigr) = 0.0336/2.7976 = 0.0120$.
Keduanya memberikan 1.2 persen.

\textbf{18.} $f = 0.75\times(-0.00019)/\bigl(40\times(-0.00147) + 0.75\times(-0.00019)\bigr)
= 0.0024$, yang meramalkan Sr \qty{10.46}{ppm} dan $^{87}\mathrm{Sr}/^{86}\mathrm{Sr} =
0.70361$, dibandingkan 0.70290 yang terukur. Asha tidak terletak pada kurva pencampuran mantel--sedimen:
sumbernya memuat komponen lain.

\textbf{19.} Kerak samudra tua yang didaur ulang, tipe HIMU: isotop Nd-nya memerlukan Sm/Nd
yang sedikit di bawah Sm/Nd \omterm{def:b2:isotope-geochemistry:reservoirs}{mantel terkuras}, isotop Sr-nya Rb/Sr rendah, dan
timbalnya (tidak diukur di sini) U/Pb tinggi, semuanya selama satu hingga dua miliar tahun.

\textbf{20.} $e^{\lambda_{147}\cdot2\times10^9} - 1 = 0.013133$. Sedimen: $0.51141 -
0.110\times0.013133 = 0.509965$; CHUR: $0.512638 - 0.1966\times0.013133 = 0.510056$;
$\varepsilon_{\mathrm{Nd}}(2\,\mathrm{Ga}) = (0.509965/0.510056 - 1)\times10^4 = -1.8$.

\textbf{21.} $f = 0.110/0.1966 - 1 = -0.440$; sejak subduksi $\varepsilon$ berubah sebesar
$25.0\times(-0.440)\times2.0 = -22.0$, sehingga $\varepsilon(2\,\mathrm{Ga}) \approx -24.0 + 22.0
= -2.0$, dalam selisih 0.2 satuan dari nilai eksak.

\textbf{22.} $e^{\lambda_{87}\cdot2\times10^9} - 1 = 0.028338$;
$0.72000 - 0.50\times0.028338 = 0.70583$.

\textbf{23.} Pada saat subduksi sedimen itu memiliki $\varepsilon_{\mathrm{Nd}} = -1.8$, dekat dengan
$-4.7$ kerak yang ketika itu baru berumur 0.7 miliar tahun, dan jauh di atas
nilai-nilai sangat negatif kerak tua: hal ini cocok dengan erosi kerak muda, mungkin
bercampur dengan sedikit bahan yang berasal dari mantel. $^{87}\mathrm{Sr}/^{86}\mathrm{Sr}$-nya
pada waktu itu, 0.7058, juga merupakan nilai batuan kerak yang muda.

\textbf{24.} Sumber Vel, pulau yang paling diperkaya, memuat
\textbf{\qty[text-series-to-math=true]{1.2}{\percent}} sedimen yang didaur ulang.
\end{solution}
